% paper/sections/appendix_numerics_solver.tex
%
% 付録（数値解析・圧力ソルバー関連の補足証明）
%   app:capillary_cfl      — 毛管波制約の導出
%   app:sweep_thomas_coeff — 欠陥補正スウィープ法の Thomas 三重対角係数導出
%   sec:dtau_derive        — 擬似時間反復の収束率解析と最適 Δτ 導出
%   app:dc_convergence_theory — 欠陥補正法の収束理論（スペクトル解析）
%   app:fvm_face_coeff     — PPE 面係数の FVM 積分論理
%
% ※ 旧 E.2, E.3, E.5 は「欠陥補正法の理論と係数導出」として1 section に統合
% ※ app:checkerboard_mode は付録 C.3（appendix_ccd_impl_s3.tex）に統合済み
%
% 内容は各 _s*.tex に分割管理されている（重複防止のため \input で委譲）

\subsubsection{毛管波制約の導出：分散関係から CFL 条件へ}
\label{app:capillary_cfl}
% ═══════════════════════════════════════════════════════════════════════════════

\textbf{Step 1：毛管波の分散関係}

気液界面上の微小振幅波（毛管波）の線形安定解析から，
角周波数 $\omega$ と波数 $k$ の関係（分散関係）は：
\begin{equation}
  \omega^2 = \frac{\sigma\,k^3}{\rho_l + \rho_g}
  \label{eq:capillary_dispersion_app}
\end{equation}
位相速度 $c_\sigma = \omega/k = \sqrt{\sigma k/(\rho_l+\rho_g)}$，
群速度 $c_g = d\omega/dk = (3/2)c_\sigma$．

\textbf{Step 2：格子上の最大分解可能波数}

格子間隔 $\Delta x$ の離散格子で分解できる最高波数はナイキスト波数
$k_{\max} = \pi/\Delta x$ である．

\textbf{Step 3：CFL 条件の適用}

$k_{\max}$ に対応する群速度で CFL 条件 $c_g(k_{\max})\,\Delta t \le \Delta x$ を適用する：
\begin{align*}
  c_g(k_{\max}) &= \frac{3}{2}\sqrt{\frac{\sigma\pi/\Delta x}{\rho_l+\rho_g}} \\
  \Delta t &\le \frac{\Delta x}{c_g(k_{\max})}
  = \frac{2}{3}\sqrt{\frac{(\rho_l+\rho_g)\,\Delta x^3}{\sigma\pi}}
\end{align*}
厳密な係数は $2/(3\sqrt{\pi}) \approx 0.376$ であるが，
文献で広く用いられる簡略形 $1/\sqrt{2\pi} \approx 0.399$ を採用すると
式 \eqref{eq:dt_sigma} を得る（$6\%$ 大きい $\Delta t$ を許すが，
毛管波 CFL の安全係数 $C_\sigma \in (0,1]$ を
$\Delta t \le C_\sigma\Delta t_\sigma$ として掛けるため実用上問題ない）：
\[
  \Delta t_\sigma = \sqrt{\frac{(\rho_l+\rho_g)\,\Delta x^3}{2\pi\sigma}}
\]

\textbf{物理的解釈：} $\Delta t_\sigma \propto \Delta x^{3/2}$ は分散関係 $\omega \propto k^{3/2}$ に由来する．
格子を半分にすると $\Delta t_\sigma$ は $1/2^{3/2} \approx 0.354$ 倍（$\approx 2.8$ 倍制約強化）．

\medskip
\textbf{大密度比（$\rho_l \gg \rho_g$）での解釈：}
$\rho_l/\rho_g = 1000$（水/空気）では $\rho_l + \rho_g \approx \rho_l$（気相寄与は $0.1\%$ 以下）．
表面張力が\textbf{復元力}（ばね），液相 $\rho_l$ が\textbf{慣性}（質量）を担い，
気相の慣性は無視できる．密度比 1000 の系では密度比 1 の系より
$\sqrt{1000} \approx 32$ 倍大きな $\Delta t_\sigma$ が許される（制約緩和）．\qed

% ═══════════════════════════════════════════════════════════════════════════════
%% 旧 appendix_numerics_schemes_s4.tex（ALE 効果）を統合
% ═══════════════════════════════════════════════════════════════════════════════

\subsubsection{時間変化格子と ALE 効果}
\label{app:grid_ale}

本付録は §\ref{sec:grid_gen} の格子が時間発展する場合の ALE（Arbitrary Lagrangian-Eulerian）補正を説明する．
\textbf{本稿の標準経路では，界面適合格子を毎ステップ再生成し，保存的再写像と速度再投影で閉じる}．
時間変化格子では，§\ref{sec:grid_ale} の
離散幾何保存則，ALE 保存更新または同等の保存的再写像，全変数再写像，
および界面面データ再構築ゲートを全て満たす必要がある．

ALE 定式化では格子速度 $\bm{v}_\text{mesh}$ を輸送方程式に追加する必要がある：
$\partial\psi/\partial t + \nabla\cdot((\bm{u} - \bm{v}_\text{mesh})\psi) = 0$．
同じ相対速度 $\bm{u}-\bm{v}_\text{mesh}$ は運動量・密度輸送にも現れるため，
$\psi$ だけを補間する格子追従は閉じた保存系ではない．

\textbf{誤差の発生機構：}
1ステップあたりの界面移動量を $\delta = |\bm{u}|\Delta t$ とすると，
ALE 補正を省略することで $\bm{v}_\text{mesh}$ による追加対流量が無視される：
\[
  \varepsilon_{\text{ALE}} \sim |\bm{v}_\text{mesh}| \cdot |\bnabla\psi| \cdot \Delta t \sim U_\Gamma \cdot \frac{1}{\varepsilon} \cdot \Delta t
\]
ここで $|\bnabla\psi| \sim 1/(4\varepsilon)$（界面近傍プロファイル），$U_\Gamma$ は界面速度である．

\textbf{近似が有効な条件：}
\begin{center}
\PaperTableSetup
\begin{tabular}{>{\raggedright\arraybackslash}p{48mm}c>{\raggedright\arraybackslash}p{42mm}}
\hline
条件 & 移動量 $\delta/h$ & ALE 省略の妥当性 \\
\hline
緩やかな界面変形 & $< 0.1$ & 良好（誤差 $< 10\%$ of $\Delta t^2$ 項） \\
標準的な問題 & $\approx 0.5$ & 条件付きで可（格子収束確認要） \\
大密度比 ($\rho_l/\rho_g{=}1000$) の高速界面 & $> 1$ & 精度劣化の可能性大（ALE 実装推奨） \\
\hline
\end{tabular}
\end{center}

$\rho_l/\rho_g = 1000$ の場合，液相変形が速く $\delta/h > 1$ となることがある．
この場合，単純な $\psi$ 補間だけの格子更新は実質的に $O(\Delta t)$ 精度（1次）となりうるため，
標準経路では体積重み付き質量補正，速度再投影，履歴リセット，面幾何再構築を同時に行う．

\textbf{Mode 1（時間固定格子，比較経路）：}
格子速度 $\bm{v}_\text{mesh} = 0$ として格子を更新しない．ALE 誤差はゼロであり，全体時間精度 $O(\Delta t^2)$ が保持される（§\ref{sec:time_accuracy_table}）．
ただし，格子は初期水準集合値 $\phi^{(0)}$ から生成されるため，界面が初期精細化領域から大きく離れる問題（例：上昇気泡ベンチマーク）では経時的に空間解像度が失われる．
長時間・長距離移動を伴う計算では，Mode 2 の界面追随格子を標準経路とする．

\textbf{Mode 2（タイムステップごとに格子を再生成，標準経路）：}
上表の $\delta/h > 1$ 条件に該当する場合，$O(\Delta t)$ に精度が劣化し，本稿の全体時間精度 $O(\Delta t^2)$ の主張（§\ref{sec:time_accuracy_table}）と矛盾する．
$\rho_l/\rho_g = 1000$（上昇気泡系）では特に注意が必要であり，Mode 2 では ALE 補正または保存的再写像を必須とする．
Mode 2 は $\psi$ 転写だけでは閉じない．
$u,v,p,\delta p,\rho,\mu,\kappa_{lg},\hat{\bm n}_{lg},\phi,\psi$ の保存的リマップと，
HFE 延長文脈・jump face cache の現在格子上での再構築を同時に満たす必要がある．
これらが欠ける場合，格子追従の利得より operator mismatch が支配的になりうるため，
式~\eqref{eq:ale_grid_gcl}--\eqref{eq:regrid_face_geometry} を標準経路の成立条件とする．

% ═══════════════════════════════════════════════════════════════════════════════
  %% E.1 毛管波制約：分散関係から CFL 条件へ（§5）

%% E.2 欠陥補正法の理論と係数導出（旧 E.2 + E.3 + E.5 を統合）
\section{欠陥補正法の理論と係数導出}
\label{app:dc_theory}
\section{スウィープ型擬似時間反復の収束率解析と最適 \texorpdfstring{$\Delta\tau$}{Δτ} の導出}
\label{sec:dtau_derive}
% ═══════════════════════════════════════════════════════════════════════════════

本付録は §\ref{sec:ppe_pseudotime} で使用した収束率指標
$\gamma(t) = (1+t^2)/(1+t)^2$ を導出し，最適擬似時間刻み $\Delta\tau_{\mathrm{opt}}$ を
理論的に確定する．

\subsection*{設定}

スウィープ型擬似時間 PPE 反復の1ステップを
\begin{equation}
  (\mathbf{I} + \Delta\tau \mathcal{L}_x)(\mathbf{I} + \Delta\tau \mathcal{L}_y)\,
  p^{m+1} = p^m + \Delta\tau\,q_h
  \label{eq:dtau_sweep}
\end{equation}
と書く．ここで $\mathcal{L}_h = \mathcal{L}_x + \mathcal{L}_y$ は
CCD 離散 Laplacian の $x$・$y$ 方向分割であり，
それぞれ固有値 $\lambda_x,\lambda_y \ge 0$ を持つ（正定値）．

\textbf{スプリッティング不動点．}
式 \eqref{eq:dtau_sweep} の不動点 $\tilde{p}$（$p^{m+1}=p^m=\tilde{p}$）は
\[
  \bigl(\mathcal{L}_h + \Delta\tau\,\mathcal{L}_x\mathcal{L}_y\bigr)\tilde{p} = q_h
\]
を満たす．これは元の PPE $\mathcal{L}_h p^* = q_h$ とは $\mathcal{O}(\Delta\tau^2)$ だけ異なり，
この差がスプリッティング誤差（\ref{warn:ppe_splitting}）の根源である．

\subsection*{誤差収縮の解析}

誤差 $e^m \equiv p^m - \tilde{p}$ を定義する．
式 \eqref{eq:dtau_sweep} から不動点方程式を辺々引くと
\[
  (\mathbf{I}+\Delta\tau\mathcal{L}_x)(\mathbf{I}+\Delta\tau\mathcal{L}_y)\,e^{m+1} = e^m
\]
となり，誤差は\textbf{斉次}の縮小方程式に従う．

\subsection*{スカラー化}

等方格子 $(\Delta x = \Delta y = h)$ かつ一様係数を仮定すると
$\lambda_x = \lambda_y = \lambda/2$（$\lambda$ は $\mathcal{L}_h$ の固有値）．
$t \equiv \Delta\tau\lambda/2$ とおけば，スペクトル解析から

\begin{equation}
  e^{m+1} = \frac{e^m}{(1+t)^2}
  \label{eq:error_contraction}
\end{equation}

が得られる．すなわちスウィープ型反復は，スプリッティング不動点 $\tilde{p}$ への誤差を
各反復ごとに $1/(1+t)^2$ 倍に縮小する（\textbf{無条件安定}；任意の $t > 0$ で縮小率 $< 1$）．

\medskip
\textbf{実用的な収束率指標 $\gamma(t)$ の導出．}
上記の純粋な誤差縮小率 $1/(1+t)^2$ は $t \to \infty$ で $0$ に近づくため，
「大きな $\Delta\tau$ が常に良い」と誤解される恐れがある．
しかし大きな $\Delta\tau$ はスプリッティング不動点誤差 $\|\tilde{p} - p^*\| \sim \mathcal{O}(\Delta\tau^2)$
をも拡大する．このトレードオフを捉える指標として，
等方格子・一様係数・スカラー固有値 $\lambda$（$t = \Delta\tau\lambda/2$）の設定で
$p^0 = 0$ から1スウィープ後の真の残差比を厳密に計算する．

\textbf{1スウィープ後の解：} 式 \eqref{eq:dtau_sweep} に $p^0 = 0$ を代入すると
$(1+t)^2 p^1 = \Delta\tau\,q_h$，すなわち $p^1 = \Delta\tau q_h/(1+t)^2$．
真の解は $\lambda p^* = q_h$，$\Delta\tau = 2t/\lambda$ より
$p^* = q_h/\lambda = \Delta\tau q_h/(2t)$．

\textbf{残差比の計算：}
\begin{align*}
  r^0 &= \lambda(p^0 - p^*) = -q_h \\
  r^1 &= \lambda(p^1 - p^*) = \lambda\!\left[\frac{\Delta\tau q_h}{(1+t)^2} - \frac{\Delta\tau q_h}{2t}\right]
       = \lambda\,\Delta\tau\,q_h \cdot \frac{2t - (1+t)^2}{2t(1+t)^2}
       = -q_h\cdot\frac{1+t^2}{(1+t)^2}
\end{align*}
ここで $2t-(1+t)^2 = -(1+t^2)$，$\lambda\,\Delta\tau = 2t$ を用いた．したがって：

\begin{equation}
  \gamma(t) \equiv \frac{|r^1|}{|r^0|} = \frac{1 + t^2}{(1+t)^2}
  \label{eq:gamma_t}
\end{equation}

（$t > 0$ に対して $\gamma(t) < 1$；初期残差を1スウィープで $\gamma$ 倍に縮小する指標．）
この指標は「収束速度」と「スプリッティング誤差下限」の両方を反映しており，
$t \to 0$ でも $t \to \infty$ でも $\gamma \to 1$（収束が遅い）という正しい定性的挙動を示す．

\subsection*{最適点の導出}

$\gamma(t)$ を最小化する $t^*$ を求める：
\begin{equation}
  \frac{d\gamma}{dt} = \frac{2(t-1)}{(1+t)^3} = 0
  \quad\Longrightarrow\quad t^* = 1
\end{equation}

このとき $\gamma(1) = 2/4 = 1/2$．すなわち最適擬似時間刻みでは
$\gamma$ の意味で\textbf{1スウィープごとに収束率指標が最小値}をとり，
収束と分割誤差のトレードオフが最も良くバランスする．

$t^* = 1$ を $\Delta\tau$ に戻すと $\Delta\tau_{\mathrm{opt}} = 2/\lambda_{\max}$．
CCD Laplacian の最大固有値は $\lambda_{\max} \approx 3.43\,a_{\max}/h^2$（$a_{\max}$ は最大逆密度係数）であるから

\begin{equation}
  \Delta\tau_{\mathrm{opt}} \approx \frac{2h^2}{3.43\,a_{\max}}
                             \approx \frac{0.58\,h^2}{a_{\max}}
  \label{eq:dtau_opt_derive}
\end{equation}

これは §\ref{sec:ppe_pseudotime} の経験的指針 $C \approx 1$ の理論的根拠となる．

\subsection*{$t \to 0$ および $t \to \infty$ の挙動}

\begin{itemize}
  \item $\gamma(t) \to 1$（$t \to 0$）：$\Delta\tau$ が小さすぎると誤差縮小率 $1/(1+t)^2 \to 1$ となり，
        スプリッティング不動点 $\tilde{p}$ への収束が遅い．
  \item $\gamma(t) \to 1$（$t \to \infty$）：$\Delta\tau$ が大きすぎるとスプリッティング不動点誤差
        $\|\tilde{p} - p^*\| \sim \mathcal{O}(\Delta\tau^2)$ が拡大し，真の解 $p^*$ からの乖離が増す．
  \item 最適点 $t=1$（$\Delta\tau = \Delta\tau_{\mathrm{opt}}$）：指標 $\gamma$ が最小値 $1/2$ をとり，
        収束速度と分割誤差のトレードオフが最良にバランスする．
        式~\eqref{eq:error_contraction} の純粋な誤差縮小率はこの点で $1/4$ であり，
        収束は速いが分割誤差下限が残る（残差収束判定により実用上は問題ない；\ref{warn:ppe_splitting} 参照）．
\end{itemize}
\qed

% ═══════════════════════════════════════════════════════════════════════════════
  %% E.2.1 Thomas 三重対角係数導出（§8）
\subsection{収束率解析と最適 \texorpdfstring{$\Delta\tau$}{Δτ} の導出}
\label{sec:dtau_derive}
% ═══════════════════════════════════════════════════════════════════════════════

本付録は §\ref{sec:ppe_pseudotime} で使用した収束率指標
$\gamma(t) = (1+t^2)/(1+t)^2$ を導出し，最適擬似時間刻み $\Delta\tau_{\mathrm{opt}}$ を
理論的に確定する．

\subsubsection*{設定}

スウィープ型擬似時間 PPE 反復の1ステップを
\begin{equation}
  (\mathbf{I} + \Delta\tau \mathcal{L}_x)(\mathbf{I} + \Delta\tau \mathcal{L}_y)\,
  p^{m+1} = p^m + \Delta\tau\,q_h
  \label{eq:dtau_sweep}
\end{equation}
と書く．ここで $\mathcal{L}_h = \mathcal{L}_x + \mathcal{L}_y$ は
PPE 演算子 $-\nabla\cdot(1/\rho\,\nabla)$ の CCD 離散化の $x$・$y$ 方向分割である．
符号規約に注意：$\mathcal{L}_h \equiv -\mathbf{L}_{\mathrm{CCD}}^{\rho}$
（付録~\ref{app:ccd_kronecker} の $\mathbf{L}_{\mathrm{CCD}}^{\rho}$ は
$\nabla\cdot(1/\rho\,\nabla)$ の離散化で負半定値；
本付録の $\mathcal{L}_h$ はその符号反転で固有値 $\lambda_x,\lambda_y \ge 0$）．
この符号規約により式~\eqref{eq:pseudo_parabolic} の誤差方程式
$\partial e/\partial\tau = -\mathcal{L}_h\,e$ が安定（減衰）となる．

\textbf{スプリッティング不動点．}
式 \eqref{eq:dtau_sweep} の不動点 $\tilde{p}$（$p^{m+1}=p^m=\tilde{p}$）は
\[
  \bigl(\mathcal{L}_h + \Delta\tau\,\mathcal{L}_x\mathcal{L}_y\bigr)\tilde{p} = q_h
\]
を満たす．これは元の PPE $\mathcal{L}_h p^* = q_h$ とは
演算子レベルで $\Delta\tau\,\mathcal{L}_x\mathcal{L}_y$ だけ異なり（解の差は $\mathcal{O}(\Delta\tau)$），
この差がスプリッティング誤差（~\ref{warn:ppe_splitting}）の根源である．

\subsubsection*{誤差収縮の解析}

誤差 $e^m \equiv p^m - \tilde{p}$ を定義する．
式 \eqref{eq:dtau_sweep} から不動点方程式を辺々引くと
\[
  (\mathbf{I}+\Delta\tau\mathcal{L}_x)(\mathbf{I}+\Delta\tau\mathcal{L}_y)\,e^{m+1} = e^m
\]
となり，誤差は\textbf{斉次}の縮小方程式に従う．

\subsubsection*{スカラー化}

等方格子 $(\Delta x = \Delta y = h)$ かつ一様係数を仮定すると
$\lambda_x = \lambda_y = \lambda/2$（$\lambda$ は $\mathcal{L}_h$ の固有値）．
$t \equiv \Delta\tau\lambda/2$ とおけば，スペクトル解析から

\begin{equation}
  e^{m+1} = \frac{e^m}{(1+t)^2}
  \label{eq:error_contraction}
\end{equation}

が得られる．すなわちスウィープ型反復は，スプリッティング不動点 $\tilde{p}$ への誤差を
各反復ごとに $1/(1+t)^2$ 倍に縮小する（\textbf{無条件安定}；任意の $t > 0$ で縮小率 $< 1$）．

\medskip
\textbf{実用的な収束率指標 $\gamma(t)$ の導出．}
上記の純粋な誤差縮小率 $1/(1+t)^2$ は $t \to \infty$ で $0$ に近づくため，
「大きな $\Delta\tau$ が常に良い」と誤解される恐れがある．
しかし大きな $\Delta\tau$ はスプリッティング不動点誤差 $\|\tilde{p} - p^*\| \sim \mathcal{O}(\Delta\tau^2)$
をも拡大する．このトレードオフを捉える指標として，
等方格子・一様係数・スカラー固有値 $\lambda$（$t = \Delta\tau\lambda/2$）の設定で
$p^0 = 0$ から1スウィープ後の真の残差比を厳密に計算する．

\textbf{1スウィープ後の解：} 式 \eqref{eq:dtau_sweep} に $p^0 = 0$ を代入すると
$(1+t)^2 p^1 = \Delta\tau\,q_h$，すなわち $p^1 = \Delta\tau q_h/(1+t)^2$．
真の解は $\lambda p^* = q_h$，$\Delta\tau = 2t/\lambda$ より
$p^* = q_h/\lambda = \Delta\tau q_h/(2t)$．

\textbf{残差比の計算：}
\begin{align*}
  r^0 &= \lambda(p^0 - p^*) = -q_h \\
  r^1 &= \lambda(p^1 - p^*) = \lambda\!\left[\frac{\Delta\tau q_h}{(1+t)^2} - \frac{\Delta\tau q_h}{2t}\right]
       = \lambda\,\Delta\tau\,q_h \cdot \frac{2t - (1+t)^2}{2t(1+t)^2}
       = -q_h\cdot\frac{1+t^2}{(1+t)^2}
\end{align*}
ここで $2t-(1+t)^2 = -(1+t^2)$，$\lambda\,\Delta\tau = 2t$ を用いた．したがって：

\begin{equation}
  \gamma(t) \equiv \frac{|r^1|}{|r^0|} = \frac{1 + t^2}{(1+t)^2}
  \label{eq:gamma_t}
\end{equation}

（$t > 0$ に対して $\gamma(t) < 1$；初期残差を1スウィープで $\gamma$ 倍に縮小する指標．）
この指標は「収束速度」と「スプリッティング誤差下限」の両方を反映しており，
$t \to 0$ でも $t \to \infty$ でも $\gamma \to 1$（収束が遅い）という正しい定性的挙動を示す．

\subsubsection*{最適点の導出}

$\gamma(t)$ を最小化する $t^*$ を求める：
\begin{equation}
  \frac{d\gamma}{dt} = \frac{2(t-1)}{(1+t)^3} = 0
  \quad\Longrightarrow\quad t^* = 1
\end{equation}

このとき $\gamma(1) = 2/4 = 1/2$．すなわち最適擬似時間刻みでは
$\gamma$ の意味で\textbf{1スウィープごとに収束率指標が最小値}をとり，
収束と分割誤差のトレードオフが最も良くバランスする．

$t^* = 1$ を $\Delta\tau$ に戻すと $\Delta\tau_{\mathrm{opt}} = 2/\lambda_{\max}$．
CCD Laplacian の最大固有値は $\lambda_{\max} = 9.6\,a_{\max}/h^2$（$a_{\max}$ は最大逆密度係数）であるから

\begin{equation}
  \Delta\tau_{\mathrm{opt}} \approx \frac{2h^2}{9.6\,a_{\max}}
                             \approx \frac{0.208\,h^2}{a_{\max}}
  \label{eq:dtau_opt_derive}
\end{equation}

これは §\ref{sec:ppe_pseudotime} の経験的指針 $C \approx 1$ の理論的根拠となる．

\subsubsection*{\texorpdfstring{$t \to 0$}{t → 0} および \texorpdfstring{$t \to \infty$}{t → ∞} の挙動}

\begin{itemize}
  \item $\gamma(t) \to 1$（$t \to 0$）：$\Delta\tau$ が小さすぎると誤差縮小率 $1/(1+t)^2 \to 1$ となり，
        スプリッティング不動点 $\tilde{p}$ への収束が遅い．
  \item $\gamma(t) \to 1$（$t \to \infty$）：$\Delta\tau$ が大きすぎるとスプリッティング不動点誤差
        $\|\tilde{p} - p^*\| \sim \mathcal{O}(\Delta\tau^2)$ が拡大し，真の解 $p^*$ からの乖離が増す．
  \item 最適点 $t=1$（$\Delta\tau = \Delta\tau_{\mathrm{opt}}$）：指標 $\gamma$ が最小値 $1/2$ をとり，
        収束速度と分割誤差のトレードオフが最良にバランスする．
        式~\eqref{eq:error_contraction} の純粋な誤差縮小率はこの点で $1/4$ であり，
        収束は速いが分割誤差下限が残る（残差収束判定により実用上は問題ない；~\ref{warn:ppe_splitting} 参照）．
\end{itemize}
\qed

\subsubsection*{CCD Laplacian の最大固有値 \texorpdfstring{$\lambda_{\max} = 9.6\,a_{\max}/h^2$}{λ\_max = 9.6 a\_max/h²} の導出}
\label{sec:lambda_max_derive}

§\ref{sec:ppe_pseudotime} の $\Delta\tau_{\mathrm{opt}} \approx 0.208\,h^2/a_{\max}$ は
$\lambda_{\max} = 2/\Delta\tau_{\mathrm{opt}} = 9.6\,a_{\max}/h^2$ を前提とする．
ここでは，この値の理論的な上限と，実際の数値的由来を示す．

\medskip
\noindent\textbf{Eq-II 単独の Fourier 解析（分離近似）．}

CCD 第2微分方程式（Eq-II）を Fourier モード $f_j = e^{\mathrm{i}\xi j}$，
$f''_j = -\Lambda/h^2 \cdot e^{\mathrm{i}\xi j}$（$\Lambda = \lambda^* h^2 > 0$）で解析すると，
\[
  \beta_2\!\left(-\frac{\Lambda}{h^2}\right)\!(e^{\mathrm{i}\xi}+e^{-\mathrm{i}\xi})
  - \frac{\Lambda}{h^2}
  = \frac{a_2}{h^2}(e^{\mathrm{i}\xi}+e^{-\mathrm{i}\xi}-2)
\]
整理して
\begin{equation}
  \Lambda(\xi) = \frac{4a_2 \sin^2(\xi/2)}{1 + 2\beta_2 \cos\xi}
  \label{eq:lambda_fourier}
\end{equation}
$\beta_2=-1/8$，$a_2=3$ を代入し，$\cos\xi = 1-2\sin^2(\xi/2)$ を用いると
$\Lambda(\xi) = 24\sin^2(\xi/2)/(3/2 + \sin^2(\xi/2))$．
$d\Lambda/d(\sin^2(\xi/2)) > 0$（単調増加）より最大値は $\xi=\pi$（$\sin^2(\pi/2)=1$）で達成される：
\begin{equation}
  \Lambda_{\max}^{\mathrm{Eq\text{-}II}} = \frac{24\cdot 1}{3/2 + 1} = \frac{24}{5/2} = 9.6
  \quad\Longrightarrow\quad
  \lambda_{\max}^{\mathrm{Eq\text{-}II}} = \frac{9.6}{h^2}
  \label{eq:lambda_max_theoretical}
\end{equation}

\medskip
\noindent\textbf{Eq-I・Eq-II 連立系の効果と数値的最大固有値．}

CCD は Eq-I（第1微分）と Eq-II（第2微分）を\textbf{同時に}解く連立系であり，
Eq-II の右辺に $f'$ 値との結合項（係数 $b_2=3/2$）が含まれる．
この結合項は $\xi=\pi$（$\sin\pi=0$）では消滅するため，
$\xi=\pi$ での最大値 $\Lambda(\pi)=9.6$ は連立系でも変わらない．
しかし，Neumann 境界条件（$f'_0=0$，$f'_{N-1}=0$ を境界行で強制）を持つ
有限格子数 $N$ の境界値問題では，周期 Fourier 解析の固有値と実際の行列スペクトルが異なる．

§~\ref{sec:verify_ccd_spectral_radius} の数値計測により，
CCD $D^{(2)}$ 行列の最大固有値は壁境界条件の有限格子（$N = 32$--$256$）でも
$\lambda_{\max} h^2 = 9.60 \pm 0.01$ であり，
Fourier 解析の理論値~\eqref{eq:lambda_max_theoretical} と一致する．
以前の $\approx 3.43/h^2$ は PPE 擬似時間ステップの実験的チューニングに
由来する値であり，$D^{(2)}$ のスペクトル半径ではない．

\begin{tcolorbox}[mybox, title={CCD $D^{(2)}$ スペクトル半径のまとめ}]
\begin{center}
\PaperTableSetup
\begin{tabular}{lc}
\hline
解析の種類 & $\lambda_{\max} h^2$（$a=1$）\\
\hline
Eq-II 単独 Fourier 解析（式~\eqref{eq:lambda_max_theoretical}） & $9.6$ \\
Eq-I+II 連立 Fourier 解析（$\xi=\pi$，coupling 消滅） & $9.6$ \\
壁 BC 有限格子（$N = 32$--$256$，数値固有値；§~\ref{sec:verify_ccd_spectral_radius}） & $9.60 \pm 0.01$ \\
FD 中心差分 $D^{(2)}_{\mathrm{FD}}$（同一格子） & $4.00$ \\
\hline
\end{tabular}
\end{center}
Fourier 理論値 $9.6/h^2$ は有限格子でも維持される（FD 中心差分の 2.4 倍）．
明示的粘性 CFL 安定限界は $\nu\Delta t/h^2 < 1/(2 \times 9.6) \approx 0.052$ であり，
CN 粘性スキーム（§~\ref{sec:verify_cn_viscous}）によりこの制約は回避される．
\end{tcolorbox}

% ═══════════════════════════════════════════════════════════════════════════════
  %% E.2.2 収束率解析と最適 Δτ（§8）
% paper/sections/appendix_numerics_solver_s5.tex
%
% 付録 D.1.3：欠陥補正法の収束理論
%   — 反復行列 M = I - ω L_L^{-1} L_H のスペクトル解析
%   — 変密度場における固有値分布
%   — 緩和係数 ω の理論的導出
%   — §8.3 (sec:defect_correction_main) の理論的裏付け

\subsection{収束理論：反復行列のスペクトル解析}
\label{app:dc_convergence_theory}

本付録は §\ref{sec:defect_correction_main} で定義した欠陥補正法（式~\eqref{eq:dc_three_step}）の
収束性を反復行列のスペクトル解析に基づいて理論的に保証する．

\subsubsection*{反復行列のスペクトル半径}

欠陥補正法の更新式（式~\eqref{eq:dc_three_step}）を誤差 $e^{(k)} \equiv p^{(k)} - p^*$
（$p^* = L_H^{-1} b$ は真の解）について整理すると：
\begin{equation}
  e^{(k+1)} = \underbrace{\bigl(I - \omega\,L_L^{-1}\,L_H\bigr)}_{\displaystyle \equiv M}\,e^{(k)}
  \label{eq:dc_iteration_matrix}
\end{equation}
反復の収束条件はスペクトル半径 $\rho(M) < 1$ であり，
$k$ 回反復後の誤差は $\|e^{(k)}\| \leq \rho(M)^k\,\|e^{(0)}\|$ で上界される．

\subsubsection*{\texorpdfstring{$M$ の固有値構造}{M の固有値構造}}

$L_L$ と $L_H$ が同一の固有ベクトル系 $\{v_j\}$ を共有すると仮定する
（等方一様格子で近似的に成立）．
$L_L\,v_j = \mu_j\,v_j$，$L_H\,v_j = \lambda_j\,v_j$ とすると
$M$ の固有値は：
\begin{equation}
  \sigma_j = 1 - \omega\,\frac{\lambda_j}{\mu_j}
  \label{eq:dc_eigenvalue}
\end{equation}
ここで $\lambda_j / \mu_j$ は CCD 演算子と FD 演算子の固有値比であり，
低周波モード（$\xi \to 0$）では $\lambda_j / \mu_j \to 1$（両者とも $\Ord{\xi^2}$ で一致），
高周波モード（$\xi \to \pi$）では $\lambda_j / \mu_j \neq 1$ となる（高次精度項の差異）．

\subsubsection*{固有値比の解析（完全連立 Eq-I+II 系）}

CCD は Eq-I（$p'$ 結合）と Eq-II（$p''$ 結合）を連立して解く．
$L_H$ の有効修正波数は両方程式を同時に解いた数値的固有値として定まり，
Eq-II 単体の近似式（$\Lambda_H \propto 24\sin^2(\xi/2)$）とは異なる．

滑らかモード（$\xi \to 0$）では CCD も FD もともに2階微分の厳密値に漸近するため
$\lambda_j/\mu_j \to 1$ となる．Nyquist モード（$\xi = \pi$）では完全連立解析
（Eq-I + Eq-II の行列固有値問題）により $\lambda_j/\mu_j = 2.4$ が得られる
（均一密度1次元，$N=32$ での数値検証と一致）．

\begin{center}
\PaperTableSetup
\begin{tabular}{ccccc}
\toprule
$kh/\pi$ & $\lambda_H h^2$ & $\lambda_{FD} h^2$ & 比 $r(k)$ & $\omega_{\max}$ \\
\midrule
$0\phantom{.00}$ (smooth) & $\approx -k^2 h^2$ & $\approx -k^2 h^2$ & $1.000$ & $2.000$ \\
$0.50$ & $-4.957$ & $-4.000$ & $1.239$ & $1.614$ \\
$0.75$ & $-11.660$ & $-6.828$ & $1.707$ & $1.171$ \\
$0.90$ & $-17.269$ & $-7.804$ & $2.213$ & $0.904$ \\
$1.00$ (Nyquist) & $-19.200$ & $-8.000$ & $2.400$ & $0.833$ \\
\bottomrule
\end{tabular}
\end{center}

固有値比は $\xi$ とともに\textbf{単調増加}し，
$[\alpha_{\min}, \alpha_{\max}] = [1.0, 2.4]$ の範囲をとる
（Nyquist が最大，滑らかモードが最小）．
\label{eq:dc_eigenvalue_ratio}

\subsubsection*{\texorpdfstring{最適緩和係数 $\omega^*$ の導出}{最適緩和係数 ω* の導出}}

$|\sigma_j| = |1 - \omega\,\lambda_j/\mu_j| < 1$ が全モードで成立するために，
$\omega < 2 / \alpha_{\max}$ が必要十分条件である：
\begin{equation}
  \omega < \frac{2}{\alpha_{\max}} = \frac{2}{2.4} \approx 0.833
  \label{eq:dc_omega_max}
\end{equation}
$\rho(M)$ を最小化する最適 $\omega^*$ は Chebyshev 中点公式より：
\begin{equation}
  \omega^* = \frac{2}{\alpha_{\min} + \alpha_{\max}}
  = \frac{2}{1.0 + 2.4}
  = \frac{2}{3.4}
  \approx 0.588
  \label{eq:dc_omega_opt}
\end{equation}
このとき
\begin{equation}
  \rho(M^*) = \frac{\alpha_{\max} - \alpha_{\min}}{\alpha_{\max} + \alpha_{\min}}
  = \frac{2.4 - 1.0}{2.4 + 1.0}
  = \frac{1.4}{3.4}
  \approx 0.41
  \label{eq:dc_rho_opt}
\end{equation}

実用的な $\omega$ 選択と各モードの減衰率を以下に示す：

\begin{center}
\PaperTableSetup
\begin{tabular}{cccc}
\toprule
$\omega$ & Nyquist $|\mu|$ & 滑らか $|\mu|$ & 評価 \\
\midrule
$1.0$（緩和なし） & $1.40$ → 発散 & $0.00$ & 不可 \\
$0.833$（$\omega_{\max}$） & $1.00$ → 停滞 & $0.17$ & 条件付 \\
$0.70$ & $0.28$ & $0.30$ & 合格（速い） \\
$\mathbf{0.50}$ & $\mathbf{0.20}$ & $\mathbf{0.50}$ & 合格（バランス型） \\
$0.30$ & $0.28$ & $0.70$ & 合格（安全側） \\
\bottomrule
\end{tabular}
\end{center}

exp10\_18（§\ref{sec:verify_dc_spectral}）では $\omega \in \{0.3, 0.5, 0.7\}$ すべてで
等密度（$\rho=1$）条件での収束を実験的に確認した（22--72 反復，tol=$10^{-8}$）．
$\omega = 0.83$ は Nyquist 減衰率 $0.992 \approx 1$ となり停滞する．
なお，$k = 3$ 反復で $\Ord{h^6}$ 精度に到達する（§\ref{sec:dc_iteration_accuracy}）理由は，
直接 LU（\texttt{spsolve}）による Step~2 の厳密求解により1反復あたりの有効減衰が大きいためである．

\subsubsection*{緩和係数 $\omega$ の役割}

式~\eqref{eq:dc_three_step} の Step~3 における $\omega$ は\textbf{直接緩和係数}であり，
Step~2 で直接 LU により厳密に求解した補正量 $\delta p^{(k+1)} = L_L^{-1} d^{(k)}$ に
スカラー倍 $\omega$ を適用する（式~\eqref{eq:dc_step2_lu}）．
収束条件 $\rho(M) < 1$ は固有値比 $\alpha_j \in [1.0, 2.4]$ の上限から
$\omega < 2/2.4 = 0.833$ として直接導かれる（式~\eqref{eq:dc_omega_max}）．

\subsubsection*{変密度場への拡張}

密度比 $\rho_l/\rho_g \gg 1$ の気液二相流では，
$L_H$ と $L_L$ の固有値比 $\alpha_j = \lambda_j/\mu_j$ が空間的に変動する．
特に：
\begin{itemize}
  \item 低密度相（気体）：$1/\rho$ が大きく $L_H$ の有効固有値が増大 → $\alpha_j$ が拡大
  \item 高密度相（液体）：$1/\rho$ が小さく有効固有値が低下 → $\alpha_j$ が縮小
\end{itemize}
この空間変動は $[\alpha_{\min}, \alpha_{\max}]$ の幅を広げ，
式~\eqref{eq:dc_rho_opt} の $\rho(M^*)$ を増大させる（収束が遅くなる）．
さらに，界面近傍では CCD の $\nabla\rho$ 評価（$\Ord{h^6}$）と
FD 前処理の $\nabla\rho$ 評価（$\Ord{h^2}$）の差異が「残差床」を形成し，
高密度比では DC+LU の収束が妨げられる（実験的確認: §\ref{sec:verify_dc_spectral}）．

\begin{tcolorbox}[resultbox, title={欠陥補正法の収束特性まとめ（DC + 直接 LU，Step 2 = \texttt{spsolve}）}]
\begin{center}
\PaperTableSetup
\begin{tabular}{lcc}
\toprule
条件 & $\rho_{\mathrm{Nyq}}(M)$ & 収束判定（tol $= 10^{-8}$） \\
\midrule
等密度，$\omega = 0.50$ & $0.20$ & 38 反復（$N=32$）収束 \\
等密度，$\omega = 0.70$ & $0.28$ & 22 反復（$N=32$）収束 \\
等密度，$\omega = 0.83\,(\omega_{\max})$ & $0.99$ & 停滞（不可） \\
変密度（$\rho_l/\rho_g = 2$，$N=64$） & $\approx 0.20$--$0.28$ & 39--74 反復 収束 \\
変密度（$\rho_l/\rho_g \ge 5$，$N=64$） & — & 界面 CCD/FD 不整合で停滞 \\
\bottomrule
\end{tabular}
\end{center}
収束には $\omega \in (0,\, 2/r_{\max}) = (0,\, 0.833)$ が必要である（実験検証: §\ref{sec:verify_dc_spectral}）．
変密度・鋭い界面では CCD（$\Ord{h^6}$）と FD（$\Ord{h^2}$）の $\nabla\rho$ 離散化の差異が
残差床 $\propto (\rho_l/\rho_g) \cdot h^2$ を生じるため，DC+LU の適用範囲は低密度比に限られる．
高密度比では一括 smoothed-Heaviside PPE の DC+LU に依存せず，
§~\ref{sec:split_ppe_framework} の分相 PPE + HFE 経路，または安定な FVM direct base solver を用いる．
\end{tcolorbox}
  %% E.2.3 収束理論：反復行列のスペクトル解析（§8）

\section{\texorpdfstring{PPE 面係数 $a_f = 2/(\rho_L+\rho_R)$ の FVM 積分論理と周期境界対応}{PPE 面係数 a\_f = 2/(ρ\_L+ρ\_R) の FVM 積分論理と周期境界対応}}
\label{app:fvm_face_coeff}
% ═══════════════════════════════════════════════════════════════════════════════

PPE 演算子の係数 $a(x) = 1/\rho(x)$ について，
セル $P$（位置 $x_L$）と $E$（位置 $x_R$）の間の等価面係数を FVM 積分で求める．

フラックス $F = a_f(p_E - p_P)/\Delta x$ において，
連続体のフラックス $F = a(x)\,dp/dx$ を区間 $[x_L, x_R]$ で積分して面係数 $a_f$ を求める．
$a(x)$ が区間内で区分的定数（各セル内で $a_L = 1/\rho_L$，$a_R = 1/\rho_R$）のとき：
\[
  \int_{x_L}^{x_R} \frac{dx}{a(x)}
  = \frac{\Delta x/2}{a_L} + \frac{\Delta x/2}{a_R}
  \quad\Rightarrow\quad
  a_f = \frac{\Delta x}{\dfrac{\Delta x/2}{a_L}+\dfrac{\Delta x/2}{a_R}}
  = \frac{2a_L a_R}{a_L+a_R}
  = \frac{2}{\rho_L+\rho_R}
\]
これは直列抵抗の等価抵抗（$1/R_\mathrm{eq} = 1/R_1+1/R_2$）と同じ論理（調和平均）であり，
熱伝導率の界面補間で使われる標準公式と同一構造を持つ．\qed

\subsection{FVM PPE 行列の周期境界条件対応}
\label{sec:fvm_periodic}
%% 旧所在: §8.4（08f_ppe_bc.tex）→ 付録 E.4 へ移動（CHK-086）

壁面（ノイマン）BC に対する FVM 離散化（式~\eqref{eq:af_exact}）は，
周期境界を持つ計算領域ではそのまま適用できない．
本稿の実装では以下の3点を変更して周期 FVM 行列を組み立てる．

\paragraph{（1）折り返しフェイスの追加}
周期 BC では内部フェイスに加え，ノード $N_x-1$ とノード $0$ を結ぶ
\textbf{折り返しフェイス（wrap face）}を追加する：
\begin{equation}
  a_{N_x - 1/2} = \frac{2}{\rho_{N_x-1} + \rho_0}
  \label{eq:af_wrap}
\end{equation}

\paragraph{（2）境界半格子体積補正の省略}
周期 BC ではすべてのノードが同じ格子幅 $h$ の制御体積を持つため，
壁面 BC の $h/2$ 体積補正は行わない（$a_f / h^2$ を一様に適用）．

\paragraph{（3）ゴーストノードの同一性拘束}
インデックス $N$ ノード（物理的に $0$ の周期像）を
\textbf{同一性拘束行}で置き換える：
\begin{equation}
  p_{\mathrm{ghost}} - p_{\mathrm{src}} = 0
  \quad\Longleftrightarrow\quad
  A[\mathrm{ghost},\,\mathrm{ghost}] = 1,\;
  A[\mathrm{ghost},\,\mathrm{src}]   = -1,\;
  b[\mathrm{ghost}] = 0
  \label{eq:ghost_identity}
\end{equation}
コーナー節点は最初の軸の拘束のみを採用する（重複適用を避けるため）．

\begin{tcolorbox}[warnbox, title={周期 FVM 行列と前処理：ILU の充填不足に注意}]
  式~\eqref{eq:ghost_identity} の同一性拘束行と FVM 行が混在した疎行列に対して，
  ILU(0)（\texttt{fill\_factor=1}）を前処理に使用すると BiCGSTAB の収束が著しく遅延する．
  実装では \texttt{fill\_factor=10} で安定収束を確認した．
  直接法（\texttt{spsolve}）は収束診断に有効である．
\end{tcolorbox}

  %% E.3 PPE 面係数の FVM 積分論理（§8）
% paper/sections/appendix_ppe_pseudotime.tex
% 付録 D.2：CCD 陰解法と仮想時間ステップ反復（PPE ソルバの実装詳細）
% ※ 旧: §8.4（本文）→ 付録 D.2 に移動（Modification IV: ソルバーアルゴリズムを付録へ）
%
\section{CCD 陰解法と仮想時間ステップ反復}
\label{sec:ppe_pseudotime}

\subsection{\texorpdfstring{初期圧力場 $p^0$ の設定}{初期圧力場 p0 の設定}}
\label{sec:ppe_p0_init}
%% 旧所在: §8.1（08_pressure.tex）→ 付録 D.2 へ移動

IPC 法の Predictor は前時刻圧力 $p^n$ を用いるため，$n=0$ における $p^0$ が必要となる．
問題に応じて以下から選択する：
\begin{itemize}
  \item \textbf{静止液滴（検証問題）：} Young--Laplace 解 $p^0 = -\sigma\kappa_0 = \sigma/R$（液相内），$p^0 = 0$（気相）を解析的に与える（$\kappa_0 = -1/R < 0$）．
  \item \textbf{一般問題（高密度比を含む）：} $n=0$ でのみ PPE（式~\eqref{eq:PPE_varrho}）を $p^0 = 0$ 初期値で解き，得られた圧力場を $p^0$ として設定する（1回追加求解のコスト）．
  \item \textbf{簡易設定：} $p^0 = 0$ のまま発進する．$\Ord{\Delta t^2}$ 精度は定常後に達成され，初期数ステップでは $\Ord{\Delta t}$ 過渡誤差が生じうるが，多くの実用問題では無視できる．
\end{itemize}

\medskip
本節で解説する\textbf{仮想時間（pseudo-time）陰解法}（スウィープ型，行列フリー）と，
付録~\ref{app:ccd_lu_direct} のクロネッカー積大域行列による \textbf{LGMRES/直接 LU 法}の
2種類の実装が本稿では使用可能である：
前者（本節）は大規模格子（$N \gtrsim 128$）でのメモリ効率に優れる matrix-free 手法，
後者（付録）はスプリッティング誤差がなくデバッグ・検証用途に適した明示的行列法である
（BiCGSTAB は低次精度実装の代替として表~\ref{tab:ppe_methods} に掲載）．
この方法は FVM ベースの BiCGSTAB（§\ref{sec:fvm_face_coeff}）に対して以下の優位性を持つ：
\textbf{空間精度}として CCD 法（6次精度）を用いるため $\Ord{h^6}$ の高精度が得られ
（表~\ref{tab:ppe_methods} 脚注 $\dagger$，収束基準は \ref{result:etol_criterion} 参照），
\textbf{行列不要}で大規模疎行列を陽に組み立てず各反復にブロック Thomas 法を適用し，
残差 $\varepsilon^m < \varepsilon_\text{tol}$ 達成後は定常解が $\mathcal{L}_h p = q_h$ を
$\varepsilon_\text{tol}$ 精度で満たす\textbf{収束後の精度保証}が得られる
（スウィープ実装の分割誤差については \ref{warn:ppe_splitting} および表~\ref{tab:ppe_methods} 参照）．

\phantomsection\label{warn:ppe_splitting}
\begin{tcolorbox}[warnbox, title={2次元求解における実装方法と分割誤差について}]
2次元 PPE の線形系 $(\mathbf{I} + \Delta\tau\mathcal{L}_h)(\delta p)^{m+1} = (\delta p)^m + \Delta\tau q_h$
（$\mathcal{L}_h = \mathcal{L}_x + \mathcal{L}_y$，式~\eqref{eq:pseudo_linear}）を実際に解く方法として，
以下の2つのアプローチがある：

\textbf{アプローチ 1（スウィープ型）：} $x$ 方向・$y$ 方向の CCD ブロック Thomas 法を交互に適用する．
実際に解く方程式は $(\mathbf{I}+\Delta\tau\mathcal{L}_x)(\mathbf{I}+\Delta\tau\mathcal{L}_y)(\delta p)^{m+1} = (\delta p)^m + \Delta\tau q_h$ であり，
これは式~\eqref{eq:pseudo_linear} の原式 $(\mathbf{I}+\Delta\tau\mathcal{L}_h)(\delta p)^{m+1} = (\delta p)^m + \Delta\tau q_h$ の\textbf{近似解}を与える．
因子展開
\[
  \bigl(\mathbf{I}+\Delta\tau\mathcal{L}_x\bigr)\bigl(\mathbf{I}+\Delta\tau\mathcal{L}_y\bigr)
  = \mathbf{I} + \Delta\tau\mathcal{L}_h + \Delta\tau^2\mathcal{L}_x\mathcal{L}_y
\]
から，原式に対して $\Delta\tau^2\mathcal{L}_x\mathcal{L}_y(\delta p)^{m+1}$ の残差（分割誤差）が生じる．

\textbf{アプローチ 2（Krylov 法）：} GMRES や CG 法で連立系を直接解く．
この場合，$\mathcal{L}_h$ の積ベクトル計算に CCD を利用でき，分割誤差は生じないが，
ブロック Thomas 法の $O(N)$ コスト優位性は失われる．

\textbf{本稿の選択と実用的な解釈：}
スウィープ型を採用する場合，「分割誤差ゼロ」は成立しないが，
\textbf{反復が収束した（$\varepsilon^m < \varepsilon_\text{tol}$）時点での解は $\mathcal{L}_h(\delta p) + \Delta\tau\mathcal{L}_x\mathcal{L}_y(\delta p) = q_h$ を満たす}．
スウィープ収束解における PPE 残差のフロアは $\Delta\tau\|\mathcal{L}_x\mathcal{L}_y(\delta p)^*\|$ であるため，
$\varepsilon_\text{tol}$ はこのフロアより大きく設定すること；
フロアより小さな $\varepsilon_\text{tol}$ を指定しても収束は達成されない．
実用上は収束後の残差を確認することで，分割誤差の影響を定量的に評価できる．
\end{tcolorbox}

\subsection{仮想時間ステップ法の定式化}
\label{sec:pseudo_formulation}

離散 PPE（IPC 増分形，未知変数 $\delta p \equiv p^{n+1}-p^n$）を
\begin{equation}
  \mathcal{L}_h\,(\delta p) = q_h,
  \qquad
  q_h \equiv \frac{(\bnabla_h^{\mathrm{RC}}\cdot\bu^*)}{\Delta t}
  \label{eq:PPE_discrete}
\end{equation}
と書く（$\mathcal{L}_h$ は式 \eqref{eq:varrho_product_rule} の product rule 展開を CCD で離散化した可変係数演算子；
具体的な行列-free 評価手順は§\ref{sec:ccd_2d_full} に示す）．
この楕円型方程式を直接解く代わりに，\textbf{仮想時間} $\tau$ を導入した放物型問題
\begin{equation}
  \pd{(\delta p)}{\tau} = q_h - \mathcal{L}_h\,(\delta p)
  \label{eq:pseudo_parabolic}
\end{equation}
を考える．$\tau\to\infty$ で $\partial(\delta p)/\partial\tau \to 0$ となり，式~\eqref{eq:pseudo_parabolic} の定常解は元の PPE~\eqref{eq:PPE_discrete} の解に一致する．

\subsection{仮想時間方向の陰的離散化}
\label{sec:pseudo_implicit}

式~\eqref{eq:pseudo_parabolic} を仮想時間方向に\textbf{陰的 Euler 法}で離散化する（$m$ は仮想時間反復ステップ）：
\begin{equation}
  \frac{(\delta p)^{m+1} - (\delta p)^m}{\Delta\tau} = q_h - \mathcal{L}_h\,(\delta p)^{m+1}
  \label{eq:pseudo_implicit_raw}
\end{equation}
整理すると，各反復ステップで解くべき線形系が得られる：
\begin{equation}
  \bigl(\mathbf{I} + \Delta\tau\,\mathcal{L}_h\bigr)\,(\delta p)^{m+1}
  = (\delta p)^m + \Delta\tau\, q_h
  \label{eq:pseudo_linear}
\end{equation}

\textbf{安定性の利点：}式~\eqref{eq:pseudo_implicit_raw} の陰的離散化は，仮想時間刻み $\Delta\tau$ の大きさに関係なく\textbf{無条件安定}である．
これにより $\Delta\tau$ を任意に大きく設定でき，収束を加速することが可能である
（CFL 制約による $\Delta\tau$ の上限が存在しない）．

\subsection{反復中の物理変数の凍結}
\label{sec:pseudo_variable_freeze}

仮想時間反復ループ（$m = 0,\,1,\,2,\ldots$）の全過程において，
$\delta p \equiv p^{n+1}-p^n$ を除くすべての物理変数は
\textbf{現在の物理時間ステップ開始時点の値に完全に凍結（フリーズ）}されなければならない．
\begin{itemize}
  \item \textbf{密度場 $\rho$：} 界面捕獲法（CLS）が現在の物理ステップで与えた
        最新の密度場を用い，反復中は界面を動かさない．
  \item \textbf{予測速度 $\bu^*$：} 運動量式（移流・粘性・表面張力等）から得た
        発散非ゼロの仮速度であり，反復中に変更しない．
  \item \textbf{物理時間刻み $\Delta t$：} 現在の流体計算で進行させようとしている
        時間刻み幅であり，定数として扱う．
\end{itemize}
この凍結により，PPE は $\delta p$ のみを未知変数とする純粋な楕円型線形問題として隔離される．
残差 $\varepsilon^m < \varepsilon_\mathrm{tol}$（式~\eqref{eq:residual}）が達成されて初めて
反復ループを抜け，次の物理計算（速度射影・界面移流）へ移行する．

\subsection{CCD 法による空間離散化}
\label{sec:pseudo_ccd_discretization}

式~\eqref{eq:pseudo_linear} の左辺行列 $\mathbf{I} + \Delta\tau\,\mathcal{L}_h$（未知量 $(\delta p)^{m+1}$）を構成するにあたり，
空間演算子 $\mathcal{L}_h$ を\textbf{CCD 法}（第\ref{sec:CCD}章）で離散化する．

2次元の PPE 演算子は
\begin{equation}
  \mathcal{L}_h = \mathcal{L}_x + \mathcal{L}_y,
  \qquad
  \mathcal{L}_x \equiv \pd{}{x}\!\Bigl(\tfrac{1}{\rho}\pd{p}{x}\Bigr),\quad
  \mathcal{L}_y \equiv \pd{}{y}\!\Bigl(\tfrac{1}{\rho}\pd{p}{y}\Bigr)
  \label{eq:L_split}
\end{equation}
と書ける．各方向の一次・二次微分を CCD で求め，$\Ord{h^6}$ の精度で $\mathcal{L}_h$ を評価する．
具体的には，各格子行（$i=$const）および格子列（$j=$const）に対して
式~\eqref{eq:ccd_global} のブロック Thomas 法を適用することで，
$\partial p/\partial x,\,\partial^2 p/\partial x^2$ および $\partial p/\partial y,\,\partial^2 p/\partial y^2$ を
$\Ord{N}$ の計算量で一括計算する．

\subsection{ADI 法との比較}
\label{sec:pseudo_adi_comparison}

ADI 法（交互方向陰解法，\cite{PeacemanRachford1955}）では，演算子を $x$・$y$ 方向に交互に分割して解く：
\begin{equation}
  \Bigl(\mathbf{I} + \tfrac{\Delta\tau}{2}\mathcal{L}_x\Bigr)\widetilde{\delta p}
  = \Bigl(\mathbf{I} - \tfrac{\Delta\tau}{2}\mathcal{L}_y\Bigr)(\delta p)^m + \tfrac{\Delta\tau}{2}\, q_h,
  \quad
  \Bigl(\mathbf{I} + \tfrac{\Delta\tau}{2}\mathcal{L}_y\Bigr)(\delta p)^{m+1}
  = \Bigl(\mathbf{I} - \tfrac{\Delta\tau}{2}\mathcal{L}_x\Bigr)\widetilde{\delta p} + \tfrac{\Delta\tau}{2}\, q_h
  \label{eq:ADI}
\end{equation}
この分割は
$
  \bigl(\mathbf{I}+\tfrac{\Delta\tau}{2}\mathcal{L}_x\bigr)
  \bigl(\mathbf{I}+\tfrac{\Delta\tau}{2}\mathcal{L}_y\bigr)
  = \mathbf{I} + \Delta\tau\mathcal{L}_h + \tfrac{\Delta\tau^2}{4}\mathcal{L}_x\mathcal{L}_y
$
の関係から，$\Ord{\Delta\tau^2}$ の\textbf{分割誤差}（splitting error）を生じる．

本稿では実装の簡便さから\textbf{スウィープ型}（$x$ 方向 Thomas 法 $\to$ $y$ 方向 Thomas 法の逐次適用）を採用する：
式~\eqref{eq:pseudo_linear} の左辺行列 $(\mathbf{I}+\Delta\tau\mathcal{L}_h)$ を
$(\mathbf{I}+\Delta\tau\mathcal{L}_x)(\mathbf{I}+\Delta\tau\mathcal{L}_y)$ で近似して因子分解し，
各反復（$(\delta p)^m \to (\delta p)^{m+1}$）を $x$・$y$ 方向の逐次 Thomas 法で実行する
（分割誤差 $\Ord{\Delta\tau^2}$ が伴うが，残差収束判定により実用上の精度を保証する；
\ref{warn:ppe_splitting} および表~\ref{tab:ppe_methods} 参照）．
ADI 法との違いは Crank--Nicolson 風の対称半分割を用いない点のみであり，
収束後の解の精度は両者とも $\varepsilon_\text{tol}$ に律速される点で同等である．

この設計から得られる利点：

\begin{itemize}
  \item 収束後の解は $\mathcal{L}_h(\delta p) = q_h$ を $\varepsilon_\text{tol}$ 精度で満たす
  \item 反復回数を残差収束判定で制御でき，過剰計算を避けられる
  \item 複雑な密度場・非直交格子への拡張が容易
\end{itemize}

（表~\ref{tab:ppe_methods} 参照）．

\subsection{欠陥補正法による離散化の分離}
\label{sec:defect_correction}

\noindent\textbf{本節で解説する欠陥補正法は，}\texttt{PPESolverSweep}\textbf{（}\texttt{ppe\_solver\_type="sweep"}\textbf{）として実装済みである．}

仮想時間法の収束先は常に CCD 精度の解を与えるが，
その収束を駆動する陰的左辺行列には必ずしも高次 CCD を用いる必要がない．
この着想を定式化した枠組みが\textbf{欠陥補正法（Defect Correction）}である．
仮想時間の更新式（式~\eqref{eq:pseudo_implicit_raw}）の空間演算子を，
\textbf{右辺（残差評価）}には高次 CCD，\textbf{左辺（陰的行列）}には低次 FD と分離する：
\begin{equation}
  \frac{(\delta p)^{(m+1)}-(\delta p)^{(m)}}{\Delta\tau}
  = \Bigl[\mathcal{L}_{FD}\bigl((\delta p)^{(m+1)}-(\delta p)^{(m)}\bigr)\Bigr]
    + R_h\bigl((\delta p)^{(m)}\bigr)
  \label{eq:defect_correction_split}
\end{equation}
ここで $R_h((\delta p)^{(m)}) \equiv q_h - \mathcal{L}_h\,(\delta p)^{(m)}$ は現在の反復値 $(\delta p)^{(m)}$
を CCD 演算子 $\mathcal{L}_h$（式~\eqref{eq:PPE_discrete}）で評価した高精度残差，
$\mathcal{L}_{FD}$ は 2次精度中心差分（FD）で構築した近似演算子であり $\mathcal{L}_h \neq \mathcal{L}_{FD}$ である．
初期値は $(\delta p)^{(0)} = 0$（\ref{box:dtau_impl} 参照）．
$\Delta_{(m)} \equiv (\delta p)^{(m+1)}-(\delta p)^{(m)}$ に整理すると，各反復で解くべき線形系は
\begin{equation}
  \Bigl[\frac{1}{\Delta\tau}\mathbf{I} - \mathcal{L}_{FD}\Bigr]\,\Delta_{(m)}
  = R_h\bigl((\delta p)^{(m)}\bigr)
  \label{eq:defect_correction_linear}
\end{equation}
となる．各項の役割を整理する：
\begin{itemize}
  \item \textbf{左辺 $[\frac{1}{\Delta\tau}\mathbf{I} - \mathcal{L}_{FD}]$：}
        $\mathcal{L}_{FD} \le 0$（負半定値）より正定値かつ対角優位であり，
        三重対角行列の集まりに近似できるため ADI 法や
        Thomas 法で $\mathcal{O}(N)$ コストで求解可能．
        ここで $\mathcal{L}_{FD}$ は 2次精度中心差分演算子であり，
        高次の $\mathcal{L}_h$（CCD）とは\textbf{異なる}演算子であることに注意する．
  \item \textbf{右辺 $R_h$：} CCD 演算子 $\mathcal{L}_h$ の $\Ord{h^6}$ 精度で評価された残差であり，
        収束先の解の空間精度を決定する．
\end{itemize}

\textbf{収束後の精度保証：}
$\Delta_{(m)} \to 0$ の極限では $R_h((\delta p)^{(m)}) \to 0$，すなわち
$\mathcal{L}_h\,(\delta p) = q_h$ が $\varepsilon_\mathrm{tol}$ の精度で満たされ，
最終圧力解は完全な CCD 空間精度を持つ（表~\ref{tab:ppe_methods} 参照）．
左辺に低次 FD を用いても\textbf{収束先の精度は損なわれない}点が欠陥補正法の本質である
（右辺 $R_h$ が高精度を保証するため）．

\textbf{実装（}\texttt{PPESolverSweep}\textbf{）との対応：}
欠陥補正法（式~\eqref{eq:defect_correction_linear}）は
\texttt{PPESolverSweep}（\texttt{ppe\_solver\_type="sweep"}）として実装されており（表~\ref{tab:ppe_methods} 5行目），
本稿の標準 PPE ソルバである．
左辺は 2次精度 FD 演算子 $\mathcal{L}_{FD}$（$x$・$y$ 方向の三重対角 Thomas 法で求解），
右辺残差は CCD 演算子 $\mathcal{L}_h$（$\Ord{h^6}$）で評価する．
$x$ 方向スウィープ後に $y$ 方向スウィープを適用する因子分解
$[\frac{1}{\Delta\tau}\mathbf{I} - \mathcal{L}_{FD_x}]^{-1}[\frac{1}{\Delta\tau}\mathbf{I} - \mathcal{L}_{FD_y}]^{-1}$
で $\Delta_{(m)}$ を求め，$(\delta p)^{m+1} = (\delta p)^{m} + \Delta_{(m)}$ と更新する
（分割誤差なし；残差収束のみが精度を決定する）．
各方向の Thomas 三重対角係数の導出は付録 \ref{app:sweep_thomas_coeff} に示す．
LTS 刻み幅（式~\eqref{eq:dtau_lts}）と組み合わせることで，
高密度比流れでも気体・液体相の残差が均等な速度で減衰する（§\ref{sec:lts_dtau}）．

\begin{table}[ht]
\centering
\small
\caption{PPE 求解法の比較}
\label{tab:ppe_methods}
\renewcommand{\arraystretch}{1.3}
\begin{tabular}{lcccc}
\hline
手法 & 空間離散化次数 & 分割誤差 & 収束判定 & 実装複雑度 \\
\hline
BiCGSTAB（FVM 離散化） & $\Ord{h^2}$ & なし & 反復回数/残差 & 低 \\
ADI + CCD（軸方向分割） & $\Ord{h^6}$ & $\Ord{\Delta\tau^2}$ & 固定スウィープ数 & 中 \\
仮想時間陰解法 + CCD（スウィープ実装）& $\Ord{h^6}^{\dagger}$ & $\Ord{\Delta\tau^2}$ & 残差駆動 & 中 \\
仮想時間陰解法 + CCD（Krylov 実装） & $\Ord{h^6}$ & なし & 残差駆動 & 高 \\
欠陥補正（FD 左辺 + CCD 残差）$\ddagger$ \textbf{[本稿標準]} & $\Ord{h^6}$ & なし & 残差駆動 & 中 \\
\hline
\multicolumn{5}{l}{$^{\dagger}$ $\varepsilon_\mathrm{tol} \ll h^6$ の設定のもとでのみ達成可能；通常の収束判定では空間精度は $\varepsilon_\mathrm{tol}$ に律速される．}\\
\multicolumn{5}{l}{\phantom{$^{\dagger}$} 分割誤差は収束後の残差が支配的でなければ影響しない（\ref{warn:ppe_splitting} 参照）．}\\
\multicolumn{5}{l}{$^{\ddagger}$ 左辺に 2次精度 FD 演算子 $\mathcal{L}_{FD}$，右辺残差に CCD 演算子 $\mathcal{L}_h$ を用いる（§\ref{sec:defect_correction}，付録 \ref{app:sweep_thomas_coeff}）．}\\
\multicolumn{5}{l}{\phantom{$^{\ddagger}$} \texttt{PPESolverSweep}（\texttt{ppe\_solver\_type="sweep"}）として実装済み；}\\
\multicolumn{5}{l}{\phantom{$^{\ddagger}$} 密度適応 LTS（§\ref{sec:lts_dtau}）により高密度比流れの収束を改善する．}
\end{tabular}
\end{table}

本稿では\textbf{欠陥補正法}（表~\ref{tab:ppe_methods} 5行目）を標準 PPE ソルバとして採用する：
左辺に 2次精度 FD 三重対角（$x$ 方向 $\to$ $y$ 方向の逐次 Thomas 法），
右辺残差に CCD（$\Ord{h^6}$）を用いることで，分割誤差なしに行列不要の $\Ord{N}$/ステップを実現する
（精度保証は §\ref{sec:defect_correction}，Thomas 係数の導出は付録 \ref{app:sweep_thomas_coeff} 参照）．

\subsection{収束判定と反復アルゴリズム}
\label{sec:pseudo_convergence}

収束判定には PPE の残差ノルムを用いる：
\begin{equation}
  \varepsilon^m \equiv \bigl\|\mathcal{L}_h(\delta p)^m - q_h\bigr\|_2
  \approx \frac{\bigl\|(\delta p)^m - (\delta p)^{m-1}\bigr\|_2}{\Delta\tau}
  \label{eq:residual}
\end{equation}
右辺の近似等号は，線形系~\eqref{eq:pseudo_linear} が各反復で厳密に解かれる場合に成立する
（厳密解時: $(\delta p)^m - (\delta p)^{m-1} = \Delta\tau(q_h - \mathcal{L}_h(\delta p)^m)$，
両辺の $\ell^2$ ノルムを $\Delta\tau$ で除すと左辺に一致）．
スウィープ実装では線形系は $\Ord{\Delta\tau^2}$ の誤差を残すため，
残差計算には増分式ではなく直接 $\mathcal{L}_h(\delta p)^m - q_h$ を評価することを推奨する．
$\varepsilon^m < \varepsilon_{\mathrm{tol}}$ を満たした時点で収束とみなし，反復を終了する．

\begin{tcolorbox}[resultbox, title={$\varepsilon_{\mathrm{tol}}$ 選択の実用基準}]
\label{result:etol_criterion}
PPE の収束誤差は速度補正 $\bu^{n+1} = \bu^* - (\Delta t/\rho)\bnabla(\delta p)$ を通じて
速度誤差 $\sim \varepsilon_\mathrm{tol}/\Delta t$ を引き起こす．
全体 $\Ord{\Delta t^2}$ 時間精度を損なわないためには，
この速度誤差が $\Ord{\Delta t^2}$ より小さいことが必要，すなわち：
\begin{equation}
  \varepsilon_\mathrm{tol} \;\ll\; \Delta t^2
  \label{eq:etol_physical}
\end{equation}
$\Delta t \sim 10^{-2}$ では $\varepsilon_\mathrm{tol} \leq 10^{-6}$，
$\Delta t \sim 10^{-3}$ では $\varepsilon_\mathrm{tol} \leq 10^{-8}$ が目安となる．
\textbf{推奨：} $\varepsilon_\mathrm{tol} = \Delta t^2 \times 10^{-2}$（$\Delta t$ 依存の適応設定）
または固定値 $\varepsilon_\mathrm{tol} = 10^{-8}$（$\Delta t \sim 10^{-3}$ 以上の場合に十分）．
なお，CCD の $\Ord{h^6}$ 空間精度を線形代数レベルで発現させるには $\varepsilon_\mathrm{tol} \ll h^6$ が
理論上の必要条件だが，NS シミュレーションでは時間精度 $\Ord{\Delta t^2}$ が支配的であり
いずれの設定も十分である（表~\ref{tab:ppe_methods} 脚注 $\dagger$ 参照）．
\end{tcolorbox}

\begin{tcolorbox}[resultbox, title={最適仮想時間刻み $\Delta\tau$ の理論（詳細：付録 §\ref{sec:dtau_derive}）}]
\label{result:dtau_opt}
等方格子・一様係数の理想条件下で，収束速度と分割誤差のトレードオフを捉える指標：
\[
  \gamma(t) = \frac{1 + t^2}{(1+t)^2},
  \qquad t \equiv \frac{\Delta\tau\,\lambda}{2}
\]
$d\gamma/dt = 0$ を解くと $t^* = 1$（$\gamma_{\min} = 1/2$）．
CCD Laplacian の最大固有値 $\lambda_{\max} = 9.6\,a_{\max}/h^2$（§~\ref{sec:verify_ccd_spectral_radius}）を代入すると：
\begin{equation}
  \Delta\tau_{\mathrm{opt}}
  \approx \frac{0.208\,h^2}{a_{\max}},
  \qquad
  a_{\max} \equiv \max_{i,j} \frac{1}{\rho_{i,j}}
  \label{eq:dtau_opt}
\end{equation}
スウィープ型は任意の $\Delta\tau > 0$ で\textbf{無条件安定}；$\Delta\tau \gg \Delta\tau_{\mathrm{opt}}$ でも $\Delta\tau \ll \Delta\tau_{\mathrm{opt}}$ でも $\gamma \to 1$（収束が遅い）．
\end{tcolorbox}

\label{box:dtau_impl}
経験則：$\Delta\tau = C_\tau \cdot h^2/a_{\max}$，$C_\tau = 1 \sim 5$（式~\eqref{eq:dtau_opt}を中心とする範囲）．

\textbf{推奨出発点：} $C_\tau = 2.0$（静止液滴テスト，$N=64$，$\rho_l/\rho_g=10$ の設定で $10\sim30$ 回で $\varepsilon_\text{tol}=10^{-8}$ 達成）．大密度比（$\rho_l/\rho_g \geq 100$）では $C_\tau=1\sim2$ から試すこと．

\textbf{初期値：} $(\delta p)^0 = 0$（増分の初期推定）；収束は典型的に数回〜数十回．初期時刻や大きな圧力変動がある場合は更に多くの反復を要する．

\subsection{密度適応型局所時間刻み（Local Time Stepping）}
\label{sec:lts_dtau}

式~\eqref{eq:dtau_opt} の最適 $\Delta\tau_\mathrm{opt} \approx 0.208\,h^2/a_\mathrm{max}$ は
等方格子・一様密度の理想条件下では有効だが，気液二相流では密度が空間的に大きく変化するため，
単一のグローバル $\Delta\tau$ では以下の問題が生じる．すなわち，同式の
$a_\mathrm{max} = \max_{i,j} 1/\rho_{i,j}$ を含むため，気体相（低密度 $\rho_\mathrm{gas}$）で
最速収束となる $\Delta\tau_\mathrm{opt}^\mathrm{gas} = \mathcal{O}(\rho_\mathrm{gas}\,h^2)$ は
液体相の最適値 $\Delta\tau_\mathrm{opt}^\mathrm{liq} = \mathcal{O}(\rho_\mathrm{liq}\,h^2)$ より
密度比倍だけ小さい．グローバル一様 $\Delta\tau$ はこの両相の\textbf{最適収束域の乖離}を解決できず，
密度比が大きいほど液体相の収束が遅くなる（剛性問題）．
この気液間の物性差に起因する\textbf{剛性（Stiffness）}を緩和する手法が
\textbf{密度適応型局所時間刻み（LTS: Local Time Stepping）}である：
\begin{equation}
  \Delta\tau_{i,j} = C_\tau \cdot \frac{\rho_{i,j}\,h^2}{2}
  \label{eq:dtau_lts}
\end{equation}
$C_\tau = 1 \sim 5$ は全領域一様のグローバル係数であり，
式~\eqref{eq:dtau_opt}・\ref{box:dtau_impl} で示した推奨範囲と同一スケールである．
式~\eqref{eq:dtau_lts} の均等化のメカニズムを示す．
1次元 PPE 演算子 $\mathcal{L}_h = \partial_x(\rho^{-1}\partial_x)$ の
セル $(i,j)$ における最適 $\Delta\tau$ の大きさは $h^2/a_{i,j}$（$a_{i,j} \equiv 1/\rho_{i,j}$）に比例する
（式~\eqref{eq:dtau_opt} 参照）．LTS 刻み幅と密度流動度の積を計算すると：
\begin{equation*}
  \Delta\tau_{i,j} \cdot a_{i,j}
  = \frac{C_\tau\,\rho_{i,j}\,h^2}{2} \cdot \frac{1}{\rho_{i,j}}
  = \frac{C_\tau\,h^2}{2} \quad (\text{全セルで一定})
\end{equation*}
となり，密度依存性が完全に打ち消される．
この均等化により，気体相・液体相を問わず各セルが同一の「規格化仮想時間刻み」で更新され，
全領域の残差がほぼ均等なペースで減衰する．

\textbf{グローバル公式との整合性：}
一様密度（$\rho_{i,j} = \rho_l$）では式~\eqref{eq:dtau_lts} は
$\Delta\tau = C_\tau\,\rho_l\,h^2/2$ となる．式~\eqref{eq:dtau_opt} の $a_\mathrm{max} = 1/\rho_l$ を代入すると
$\Delta\tau_\mathrm{opt} \approx 0.208\,h^2\,\rho_l$ であり，
$C_\tau \approx 0.416$ に相当する範囲で両式は整合する（\ref{box:dtau_impl} の $C_\tau = 1 \sim 5$ の範囲内）．

\textbf{$C_\tau$ の動的増大：}
反復初期は $C_\tau$ を小さめ（$C_\tau = \mathcal{O}(1)$）に設定して安定性を確保し，
残差が十分に減少した後に段階的に増大させると定常解への到達を劇的に早める．
$\Delta\tau_{i,j}$ が陰的 Euler（式~\eqref{eq:pseudo_implicit_raw}）に基づく限り，
$C_\tau$ を大きくしても発散しない（無条件安定性；\ref{warn:ppe_splitting} 参照）．

\subsection{全周 Neumann 境界の可解性条件とゲージ固定点の対称性}
\label{sec:ppe_gauge_neumann}

\noindent\textbf{可解性条件（Neumann 境界の必要条件）．}
全境界が $\partial(\delta p)/\partial n = 0$ の場合，
PPE~\eqref{eq:PPE_discrete} が解を持つための必要十分条件は
\begin{equation}
  \int_\Omega q_h \, dV = 0
  \quad\Longleftrightarrow\quad
  \sum_{i,j}(q_h)_{i,j}\,h^2 = 0
  \label{eq:neumann_solvability}
\end{equation}
である．発散定理から
$\sum_{i,j}(q_h)_{i,j}h^2 = (1/\Delta t)\oint_{\partial\Omega}(\bu^*\cdot\hat{n})\,dS$
であり，固体壁境界 $\bu^*\cdot\hat{n}=0$ では\textbf{自動的に満たされる}（§\ref{sec:time_level1} の壁面 BC 参照）．
数値誤差による可解性条件の微小な違反が現れる場合は，
$(q_h)_{i,j} \leftarrow (q_h)_{i,j} - \overline{q_h}$（$\overline{q_h} = $ 全域平均）の
補正を RHS に適用して厳密に $\sum q_h = 0$ を回復させる．

\medskip
\noindent\textbf{ゲージ固定点の対称性：中心セル $(N/2,\,N/2)$ の根拠．}

可解性条件が満たされても，解 $\delta p$ は定数零固有モードの自由度を持つ（§\ref{sec:pinpoint}）．
ピン条件（式~\eqref{eq:pin_overwrite}）でこの自由度を除去する際，
固定点 $(i_0,j_0)$ の選択は\textbf{数値対称性の保持}に影響する．

正方形領域 $[0,L]^2$ 上の物理設定（単一気泡・静止液滴）が
$x$-$y$ 面内4倍対称（$x\leftrightarrow L-x$，$y\leftrightarrow L-y$）を持つ場合，
圧力場も同一の格子対称性を持つ必要がある．
この対称性を数値的に保持するには，
ピン条件 $p_{i_0,j_0}=0$ が変換 $i\mapsto N-1-i$ および $j\mapsto N-1-j$ に対して不変であることが必要である：
\[
  i_0 = N-1-i_0 \;\text{かつ}\; j_0 = N-1-j_0
  \;\Longrightarrow\;
  i_0 = j_0 = \bigl\lfloor N/2 \bigr\rfloor
\]
（本稿では $N$ が偶数の正方格子を前提とし，$i_0=j_0=N/2$ を使用する；
奇数 $N$ の場合は $i_0=j_0=(N-1)/2$ が対称中心だが，本稿ではこの設定は対象としない．）

\begin{center}
\renewcommand{\arraystretch}{1.3}
\begin{tabular}{lcc}
\hline
ゲージ点 & 4倍対称の不変点 & 対称性の保持 \\
\hline
$(0,0)$（コーナー） & \texttimes & 破る \\
$(N/2,\,N/2)$（中心） & \checkmark & 保持 \\
\hline
\end{tabular}
\end{center}
コーナー点 $(0,0)$ をゲージに選ぶと，対称な物理設定でも圧力ゲージの非対称性が
速度補正 $\bu^{n+1}=\bu^*-(\Delta t/\rho)\bnabla(\delta p)$ を通じて微小な流れ場非対称性を誘起し，
長時間積分で蓄積する（symmetry-breaking；§\ref{sec:pinpoint} 原則 2 参照）．
中心ピンはこの人工的な非対称性をゼロにする選択である．

       %% E.4 CCD 陰解法と仮想時間ステップ反復（実装詳細）
